% ==============================================================================
% File: elf_paper.tex
% Bài báo: Embedded Language Flows (ELF)
% Mô hình khuếch tán ngôn ngữ liên tục dựa trên Flow Matching
% ==============================================================================

\section{Embedded Language Flows (ELF)}

% Slide 0: Giới thiệu tổng quan bài báo
\begin{frame}{ELF: Tổng quan bài báo}
    \begin{itemize}
        \item \textbf{Tên bài báo:} \textit{Embedded Language Flows (ELF)}
        \item \textbf{Ý tưởng cốt lõi:} Mô hình khuếch tán ngôn ngữ liên tục dựa trên Flow Matching trong không gian nhúng biểu diễn (latent space).
        \item \textbf{Đóng góp chính:}
        \begin{itemize}
            \item Chứng minh continuous diffusion hoàn toàn vượt trội discrete diffusion trên ngôn ngữ nếu khử nhiễu hoàn toàn trong không gian nhúng liên tục.
            \item Chỉ thực hiện rời rạc hóa (discretization/decoding) tại bước cuối cùng ($t = 1$).
            \item Sử dụng mạng chia sẻ trọng số giữa denoising và decoding, tích hợp sẵn Classifier-Free Guidance (CFG) lúc huấn luyện.
        \end{itemize}
    \end{itemize}
\end{frame}

% Slide 1: Bối cảnh & Phân loại
\begin{frame}{1. ELF: Bối cảnh \& Phân loại DLMs}
    \textbf{Bối cảnh:}
    \begin{itemize}
        \item Diffusion và flow-based models là tiêu chuẩn vàng trong sinh ảnh/video, đang được mở rộng sang xử lý ngôn ngữ.
    \end{itemize}
    \vspace{0.2cm}
    \textbf{Phân loại mô hình khuếch tán ngôn ngữ (DLMs):}
    \begin{itemize}
        \item \textit{Rời rạc (Discrete):} Thao tác trực tiếp trên token (trước đây được ưu tiên do dễ tối ưu hóa).
        \item \textit{Liên tục (Continuous):} Ánh xạ token sang không gian biểu diễn liên tục để tiến hành khử nhiễu.
    \end{itemize}
    \vspace{0.2cm}
    \textbf{Giải pháp đề xuất của ELF:}
    \begin{itemize}
        \item Khử nhiễu hoàn toàn trong không gian nhúng liên tục.
        \item Định nghĩa trường vận tốc qua thời gian liên tục với Flow Matching.
        \item Dùng mạng chia sẻ trọng số, chỉ rời rạc hóa tại bước $t = 1$.
    \end{itemize}
\end{frame}

% Slide 2: Khung làm việc ELF
\begin{frame}{2. ELF: Khung làm việc (Framework)}
    \begin{itemize}
        \item \textbf{Từ rời rạc sang liên tục:} Chuỗi token $s \in \mathcal{V}^L$ được chuyển sang biểu diễn liên tục $x$ bằng bộ mã hóa T5 đông cứng.
        \item \textbf{Nội suy tuyến tính (Rectified Flows):} Tạo biến ẩn $z_t$ dọc theo đường đi:
        $$z_t = t x + (1 - t)\epsilon, \quad t \in [0, 1], \quad \epsilon \sim \mathcal{N}(0, \mathbf{I})$$
        \item \textbf{Trường vận tốc dòng (Flow Velocity):} Đạo hàm theo thời gian:
        $$v = \frac{dz}{dt} = x - \epsilon$$
        \item \textbf{Dự đoán $x$ (x-prediction):} Trực tiếp dự đoán biểu diễn sạch $x_\theta(z_t, t)$ thay vì dự đoán vận tốc $v$, giúp mạng nơ-ron hoạt động ổn định ở số chiều cực cao.
    \end{itemize}
\end{frame}

% Slide 3: Hàm mục tiêu & Rời rạc hóa
\begin{frame}{3. ELF: Hàm mục tiêu \& Rời rạc hóa}
    \begin{itemize}
        \item \textbf{Nhánh khử nhiễu (Denoising - chiếm 80\% trọng số):} Tối thiểu hóa sai số MSE của trường vận tốc dự đoán:
        $$\mathcal{L}_{\text{MSE}} = \mathbb{E}_{t, x, \epsilon} \left[ \frac{1}{(1-t)^2} \|x_\theta(z_t, t) - x\|^2 \right]$$
        \item \textbf{Nhánh giải mã (Decoding - chiếm 20\% trọng số):}
        \begin{itemize}
            \item Tại bước cuối $t = 1$, biến đổi $z_t$ thành phân phối token qua ma trận unembedding $W$.
            \item Dùng chung trọng số với mạng khử nhiễu, huấn luyện bằng Cross-Entropy:
            $$\mathcal{L}_{\text{CE}} = \mathbb{E}_{\tilde{z}} \left[ \text{CrossEntropy}(W x_\theta(\tilde{z}), s) \right]$$
        \end{itemize}
    \end{itemize}
\end{frame}

% Slide 4: Điều kiện hóa & Dẫn hướng
\begin{frame}{4. ELF: Điều kiện hóa \& Dẫn hướng (Guidance)}
    \begin{itemize}
        \item \textbf{Self-conditioning:} Thực hiện lượt truyền thứ hai điều kiện hóa trên dự đoán trung gian $\hat{x}'$ từ bước trước để tăng độ mượt của quỹ đạo.
        \item \textbf{Classifier-Free Guidance (CFG):} Dễ dàng áp dụng trong miền liên tục:
        $$v_{\text{cfg}}(z_t \mid c) = \omega v(z_t \mid c) + (1 - \omega) v(z_t \mid \emptyset)$$
        \item \textbf{Training-time CFG:} Tích hợp thẳng CFG vào hàm mục tiêu vận tốc ngay từ lúc huấn luyện $\rightarrow$ loại bỏ hoàn toàn chi phí tính toán kép khi suy luận.
        \item \textbf{Sinh có điều kiện:} Ghép prompt vào trước chuỗi mục tiêu; phần prompt được giữ nguyên không làm nhiễu.
    \end{itemize}
\end{frame}

% Slide 5: Kết quả thực nghiệm
\begin{frame}{5. ELF: Kết quả thực nghiệm}
    \textbf{Sinh không điều kiện (OpenWebText):}
    \begin{itemize}
        \item Ở 32 bước lấy mẫu, ELF-B (105M) đạt \textbf{Gen. PPL 24.08}, vượt trội so với các đối thủ rời rạc như MDLM (104.29) và Duo (62.47).
        \item \textbf{Hiệu quả dữ liệu cực cao:} Đạt đỉnh chỉ với 45B token huấn luyện (ít hơn $\sim$10 lần so với $>500$B tokens của các baseline).
    \end{itemize}
    \vspace{0.2cm}
    \textbf{Chưng cất lũy tiến (Progressive Distillation):}
    \begin{itemize}
        \item Cho phép giảm số bước lấy mẫu xuống mức cực ít (\textbf{1 đến 4 bước}) mà vẫn duy trì chất lượng sinh văn bản xuất sắc.
    \end{itemize}
    \vspace{0.2cm}
    \textbf{Sinh có điều kiện:}
    \begin{itemize}
        \item \textbf{Tóm tắt XSum:} ROUGE-1 đạt \textbf{36.0} (đánh bại các mô hình tự hồi quy cùng kích cỡ).
        \item \textbf{Dịch máy WMT14:} Đạt BLEU \textbf{26.4}.
    \end{itemize}
\end{frame}

% Slide 6: Kết luận
\begin{frame}{6. ELF: Kết luận \& Đóng góp}
    \begin{itemize}
        \item \textbf{Khẳng định:} Mô hình khuếch tán ngôn ngữ liên tục hoàn toàn có thể cạnh tranh và vượt trội hơn các phương pháp rời rạc.
        \item \textbf{Bí quyết tối ưu của ELF:}
        \begin{enumerate}
            \item Giữ toàn bộ quỹ đạo làm nhiễu trong không gian nhúng liên tục.
            \item Chỉ thực hiện rời rạc hóa ở bước thời gian cuối cùng ($t = 1$).
            \item Sử dụng mạng chia sẻ trọng số, loại bỏ bộ giải mã phụ trợ rườm rà.
        \end{enumerate}
        \item \textbf{Tác động:} Mở ra tiềm năng tạo văn bản chất lượng cao với chi phí huấn luyện thấp và tốc độ suy luận cực nhanh.
    \end{itemize}
\end{frame}
